% GENERATED FILE. Edit canonical lessons, then run npm run generate:books.
% canonical-source-hash: fnv1a64:52d6ad4c2fd4b48a
% canonical-lessons: IT-C163-meccanico, IT-C163-elettricista, IT-C163-idraulico, IT-C163-operaio, IT-C163-impiegato

\chapter{Mechanic, Electrician, Plumber, Worker, Office worker}
\label{ch:it-mechanic-electrician-plumber-worker-office-worker}
\hlchaptermodality{163}

\begin{chapteropening}
\textbf{By the end of this chapter:} \emph{I can say a mechanic, an electrician, a plumber, a worker and an office worker.}

Complete the last lesson of chapter 163: I can say a mechanic, an electrician, a plumber, a worker and an office worker.
\end{chapteropening}

\section[il meccanico]{il meccanico — a mechanic}

\label{lesson:IT-C163-meccanico}



\begin{quote}
Before the new one: say the Italian for a battery, then the Italian for a charger.
\end{quote}

\subsection*{You'll want to know: il meccanico}
\textbf{il meccanico} — \textquotedblleft{}a mechanic\textquotedblright{}.

\begin{cousinweb}
\textbf{meccanico} $\leftarrow$ Latin \emph{mechanicus}, \textquotedblleft{}of machines\textquotedblright{}, from Greek \emph{mechanikos}, from \emph{mechane}, \textquotedblleft{}a machine\textquotedblright{} — a learned borrowing. EtymDB 2.1 has no row for \emph{meccanico}; it derives Spanish \emph{mecánico} from \emph{mechanicus}.
\end{cousinweb}

\begin{grammarlens}[title={What you've built}]
One more word for this chapter.
\end{grammarlens}

\subsection*{Your turn}
\begin{itemize}
\raggedright
  \item \emph{Say it:} \emph{il meccanico}
  \item \emph{Say it:} \emph{il meccanico}, once more
  \item \emph{Say it:} say \emph{il meccanico}
  \item \emph{Recall:} say the Italian for a battery, then the Italian for a charger, then say \emph{il meccanico} again
\end{itemize}

\begin{tcolorbox}[breakable,colback=teal!4,colframe=teal!35!black,title={Before you move on}]
What does il meccanico mean? (\textquotedblleft{}A mechanic\textquotedblright{}.) Say it once more.
\end{tcolorbox}
\section[l'elettricista]{l'elettricista — an electrician}

\label{lesson:IT-C163-elettricista}



\begin{quote}
Before the new one: say the Italian for a charger, then the Italian for a mechanic.
\end{quote}

\subsection*{You'll want to know: l'elettricista}
\textbf{l'elettricista} — \textquotedblleft{}an electrician\textquotedblright{}.

\begin{cousinweb}
\textbf{elettricista} — built from the word for \textquotedblleft{}electric\textquotedblright{} with \emph{-ista}. That word goes back to Latin \emph{electrum}, \textquotedblleft{}amber\textquotedblright{}, from Greek \emph{elektron}: rubbed amber gives a spark. (EtymDB 2.1 records the Latin and Greek words.)
\end{cousinweb}

\begin{grammarlens}[title={What you've built}]
One more word for this chapter.
\end{grammarlens}

\subsection*{Your turn}
\begin{itemize}
\raggedright
  \item \emph{Say it:} \emph{l'elettricista}
  \item \emph{Say it:} \emph{l'elettricista}, once more
  \item \emph{Say it:} say \emph{l'elettricista}
  \item \emph{Recall:} say the Italian for a charger, then the Italian for a mechanic, then say \emph{l'elettricista} again
\end{itemize}

\begin{tcolorbox}[breakable,colback=teal!4,colframe=teal!35!black,title={Before you move on}]
What does l'elettricista mean? (\textquotedblleft{}An electrician\textquotedblright{}.) Say it once more.
\end{tcolorbox}
\section[l'idraulico]{l'idraulico — a plumber}

\label{lesson:IT-C163-idraulico}



\begin{quote}
Before the new one: say the Italian for a mechanic, then the Italian for an electrician.
\end{quote}

\subsection*{You'll want to know: l'idraulico}
\textbf{l'idraulico} — \textquotedblleft{}a plumber\textquotedblright{}.

\begin{cousinweb}
\textbf{idraulico} $\leftarrow$ Latin \emph{hydraulicus}, \textquotedblleft{}of the water organ\textquotedblright{}, from Greek \emph{hydraulikos} — a learned borrowing. EtymDB 2.1 has no row for \emph{idraulico}; it derives French \emph{hydraulique} from \emph{hydraulicus}.
\end{cousinweb}

\begin{grammarlens}[title={What you've built}]
One more word for this chapter.
\end{grammarlens}

\subsection*{Your turn}
\begin{itemize}
\raggedright
  \item \emph{Say it:} \emph{l'idraulico}
  \item \emph{Say it:} \emph{l'idraulico}, once more
  \item \emph{Say it:} say \emph{l'idraulico}
  \item \emph{Recall:} say the Italian for a mechanic, then the Italian for an electrician, then say \emph{l'idraulico} again
\end{itemize}

\begin{tcolorbox}[breakable,colback=teal!4,colframe=teal!35!black,title={Before you move on}]
What does l'idraulico mean? (\textquotedblleft{}A plumber\textquotedblright{}.) Say it once more.
\end{tcolorbox}
\section[l'operaio]{l'operaio — a worker}

\label{lesson:IT-C163-operaio}



\begin{quote}
Before the new one: say the Italian for an electrician, then the Italian for a plumber.
\end{quote}

\subsection*{You'll want to know: l'operaio}
\textbf{l'operaio} — \textquotedblleft{}a worker\textquotedblright{}.

\begin{cousinweb}
\textbf{operaio} $\leftarrow$ Latin \emph{operarius}, \textquotedblleft{}a labourer\textquotedblright{}, from \emph{opera}, \textquotedblleft{}work\textquotedblright{} — inherited: Latin \emph{-arius} passed down by speech became \emph{-aio}. (EtymDB 2.1)
\end{cousinweb}

\begin{grammarlens}[title={What you've built}]
One more word for this chapter.
\end{grammarlens}

\subsection*{Your turn}
\begin{itemize}
\raggedright
  \item \emph{Say it:} \emph{l'operaio}
  \item \emph{Say it:} \emph{l'operaio}, once more
  \item \emph{Say it:} say \emph{l'operaio}
  \item \emph{Recall:} say the Italian for an electrician, then the Italian for a plumber, then say \emph{l'operaio} again
\end{itemize}

\begin{tcolorbox}[breakable,colback=teal!4,colframe=teal!35!black,title={Before you move on}]
What does l'operaio mean? (\textquotedblleft{}A worker\textquotedblright{}.) Say it once more.
\end{tcolorbox}
\section[l'impiegato]{l'impiegato — an office worker}

\label{lesson:IT-C163-impiegato}



\begin{quote}
Before the new one: say the Italian for a plumber, then the Italian for a worker.
\end{quote}

\subsection*{You'll want to know: l'impiegato}
\textbf{l'impiegato} — \textquotedblleft{}an office worker\textquotedblright{}.

\begin{cousinweb}
\textbf{impiegato} is the past participle of \emph{impiegare}, \textquotedblleft{}to employ\textquotedblright{}, which $\leftarrow$ Latin \emph{implicare}, \textquotedblleft{}to enfold, to involve\textquotedblright{}. (EtymDB 2.1)
\end{cousinweb}

\begin{grammarlens}[title={What you've built}]
One more word for this chapter.
\end{grammarlens}

\subsection*{Your turn}
\begin{itemize}
\raggedright
  \item \emph{Say it:} \emph{l'impiegato}
  \item \emph{Say it:} \emph{l'impiegato}, once more
  \item \emph{Say it:} say \emph{l'impiegato}
  \item \emph{Recall:} say the Italian for a plumber, then the Italian for a worker, then say \emph{l'impiegato} again
\end{itemize}

\begin{tcolorbox}[breakable,colback=teal!4,colframe=teal!35!black,title={Before you move on}]
What does l'impiegato mean? (\textquotedblleft{}An office worker\textquotedblright{}.) Say it once more.
\end{tcolorbox}
